% GENERATED FILE. Edit canonical lessons, then run npm run generate:books.
% canonical-source-hash: fnv1a64:326632490cb7359f
% canonical-lessons: TA-C236-arunkatciyakam, TA-C236-cantu, TA-C236-turaimukam, TA-C236-kottai, TA-C236-aranmanai

\chapter{Museum, Lane, Harbour, Fort, Palace}
\label{ch:ta-museum-lane-harbour-fort-palace}
\hlchaptermodality{236}

\begin{chapteropening}
\textbf{By the end of this chapter:} \emph{I can say a museum, a lane, a harbour, a fort and a palace.}

Complete the last lesson of chapter 236: I can say a museum, a lane, a harbour, a fort and a palace.
\end{chapteropening}

\section[\texorpdfstring{\ta{அருங்காட்சியகம்} (aruṅkāṭciyakam)}{அருங்காட்சியகம் (aruṅkāṭciyakam)}]{\ta{அருங்காட்சியகம்} (aruṅkāṭciyakam) — a museum}

\label{lesson:TA-C236-arunkatciyakam}



\begin{quote}
Before the new one: say the Tamil for east, then the Tamil for west.
\end{quote}

\subsection*{You'll want to know: \ta{அருங்காட்சியகம்}}
\textbf{\ta{அருங்காட்சியகம்}} — \emph{aruṅkāṭciyakam} — \textquotedblleft{}a museum\textquotedblright{}.

\begin{cousinweb}
\textbf{\ta{அருங்காட்சியகம்}} (\emph{aruṅkāṭciyakam}) is a modern coinage from native parts: \emph{aru}, \textquotedblleft{}rare\textquotedblright{} (DEDR 221), \emph{kāṭci}, \textquotedblleft{}sight, show\textquotedblright{} (DEDR 1443), and \emph{akam}, \textquotedblleft{}house\textquotedblright{} (DEDR 7): \textquotedblleft{}the house of rare sights\textquotedblright{}.
\end{cousinweb}

\begin{grammarlens}[title={What you've built}]
One more word for this chapter.
\end{grammarlens}

\subsection*{Your turn}
\begin{itemize}
\raggedright
  \item \emph{Say it:} \emph{aruṅkāṭciyakam}
  \item \emph{Say it:} \emph{aruṅkāṭciyakam}, once more
  \item \emph{Say it:} say \emph{aruṅkāṭciyakam}
  \item \emph{Recall:} say the Tamil for east, then the Tamil for west, then say \emph{aruṅkāṭciyakam} again
\end{itemize}

\begin{tcolorbox}[breakable,colback=teal!4,colframe=teal!35!black,title={Before you move on}]
What does \ta{அருங்காட்சியகம்} mean? (\textquotedblleft{}A museum\textquotedblright{}.) Say it once more.
\end{tcolorbox}
\section[\texorpdfstring{\ta{சந்து} (cantu)}{சந்து (cantu)}]{\ta{சந்து} (cantu) — a lane, an alley}

\label{lesson:TA-C236-cantu}



\begin{quote}
Before the new one: say the Tamil for west, then the Tamil for a museum.
\end{quote}

\subsection*{You'll want to know: \ta{சந்து}}
\textbf{\ta{சந்து}} — \emph{cantu} — \textquotedblleft{}a lane, an alley\textquotedblright{}.

\begin{cousinweb}
\textbf{\ta{சந்து}} (\emph{cantu}) is usually traced to Sanskrit \emph{saṁdhi}, \textquotedblleft{}joint, junction\textquotedblright{} (CDIAL 12913): a lane is a gap between houses.
\end{cousinweb}

\begin{grammarlens}[title={What you've built}]
One more word for this chapter.
\end{grammarlens}

\subsection*{Your turn}
\begin{itemize}
\raggedright
  \item \emph{Say it:} \emph{cantu}
  \item \emph{Say it:} \emph{cantu}, once more
  \item \emph{Say it:} say \emph{cantu}
  \item \emph{Recall:} say the Tamil for west, then the Tamil for a museum, then say \emph{cantu} again
\end{itemize}

\begin{tcolorbox}[breakable,colback=teal!4,colframe=teal!35!black,title={Before you move on}]
What does \ta{சந்து} mean? (\textquotedblleft{}A lane, an alley\textquotedblright{}.) Say it once more.
\end{tcolorbox}
\section[\texorpdfstring{\ta{துறைமுகம்} (tuṟaimukam)}{துறைமுகம் (tuṟaimukam)}]{\ta{துறைமுகம்} (tuṟaimukam) — a harbour, a port}

\label{lesson:TA-C236-turaimukam}



\begin{quote}
Before the new one: say the Tamil for a museum, then the Tamil for a lane.
\end{quote}

\subsection*{You'll want to know: \ta{துறைமுகம்}}
\textbf{\ta{துறைமுகம்}} — \emph{tuṟaimukam} — \textquotedblleft{}a harbour, a port\textquotedblright{}.

\begin{cousinweb}
\textbf{\ta{துறைமுகம்}} (\emph{tuṟaimukam}) is native: \emph{tuṟai}, \textquotedblleft{}harbour, landing place\textquotedblright{} (DEDR 3370), and \emph{mukam}, \textquotedblleft{}face, front\textquotedblright{} (DEDR 4889).
\end{cousinweb}

\begin{grammarlens}[title={What you've built}]
One more word for this chapter.
\end{grammarlens}

\subsection*{Your turn}
\begin{itemize}
\raggedright
  \item \emph{Say it:} \emph{tuṟaimukam}
  \item \emph{Say it:} \emph{tuṟaimukam}, once more
  \item \emph{Say it:} say \emph{tuṟaimukam}
  \item \emph{Recall:} say the Tamil for a museum, then the Tamil for a lane, then say \emph{tuṟaimukam} again
\end{itemize}

\begin{tcolorbox}[breakable,colback=teal!4,colframe=teal!35!black,title={Before you move on}]
What does \ta{துறைமுகம்} mean? (\textquotedblleft{}A harbour, a port\textquotedblright{}.) Say it once more.
\end{tcolorbox}
\section[\texorpdfstring{\ta{கோட்டை} (kōṭṭai)}{கோட்டை (kōṭṭai)}]{\ta{கோட்டை} (kōṭṭai) — a fort}

\label{lesson:TA-C236-kottai}



\begin{quote}
Before the new one: say the Tamil for a lane, then the Tamil for a harbour.
\end{quote}

\subsection*{You'll want to know: \ta{கோட்டை}}
\textbf{\ta{கோட்டை}} — \emph{kōṭṭai} — \textquotedblleft{}a fort\textquotedblright{}.

\begin{cousinweb}
\textbf{\ta{கோட்டை}} (\emph{kōṭṭai}) is native Dravidian. DEDR 2207 lists Malayalam \emph{kōṭṭa}, \textquotedblleft{}fort\textquotedblright{} and Kannada \emph{kōṭe}. The entry compares Sanskrit \emph{koṭṭa}, \textquotedblleft{}fort\textquotedblright{}.
\end{cousinweb}

\begin{grammarlens}[title={What you've built}]
One more word for this chapter.
\end{grammarlens}

\subsection*{Your turn}
\begin{itemize}
\raggedright
  \item \emph{Say it:} \emph{kōṭṭai}
  \item \emph{Say it:} \emph{kōṭṭai}, once more
  \item \emph{Say it:} say \emph{kōṭṭai}
  \item \emph{Recall:} say the Tamil for a lane, then the Tamil for a harbour, then say \emph{kōṭṭai} again
\end{itemize}

\begin{tcolorbox}[breakable,colback=teal!4,colframe=teal!35!black,title={Before you move on}]
What does \ta{கோட்டை} mean? (\textquotedblleft{}A fort\textquotedblright{}.) Say it once more.
\end{tcolorbox}
\section[\texorpdfstring{\ta{அரண்மனை} (araṇmaṉai)}{அரண்மனை (araṇmaṉai)}]{\ta{அரண்மனை} (araṇmaṉai) — a palace}

\label{lesson:TA-C236-aranmanai}



\begin{quote}
Before the new one: say the Tamil for a harbour, then the Tamil for a fort.
\end{quote}

\subsection*{You'll want to know: \ta{அரண்மனை}}
\textbf{\ta{அரண்மனை}} — \emph{araṇmaṉai} — \textquotedblleft{}a palace\textquotedblright{}.

\begin{cousinweb}
\textbf{\ta{அரண்மனை}} (\emph{araṇmaṉai}) is \emph{araṇ}, \textquotedblleft{}defence, fortress\textquotedblright{}, which DEDR 201 traces to Sanskrit \emph{śaraṇa}, \textquotedblleft{}refuge\textquotedblright{}, and native \emph{maṉai}, \textquotedblleft{}house\textquotedblright{} (DEDR 4776).
\end{cousinweb}

\begin{grammarlens}[title={What you've built}]
One more word for this chapter.
\end{grammarlens}

\subsection*{Your turn}
\begin{itemize}
\raggedright
  \item \emph{Say it:} \emph{araṇmaṉai}
  \item \emph{Say it:} \emph{araṇmaṉai}, once more
  \item \emph{Say it:} say \emph{araṇmaṉai}
  \item \emph{Recall:} say the Tamil for a harbour, then the Tamil for a fort, then say \emph{araṇmaṉai} again
\end{itemize}

\begin{tcolorbox}[breakable,colback=teal!4,colframe=teal!35!black,title={Before you move on}]
What does \ta{அரண்மனை} mean? (\textquotedblleft{}A palace\textquotedblright{}.) Say it once more.
\end{tcolorbox}
